% Methodology — Securing LLMs Against Prompt Injection
% Academic manuscript section (IEEE / arXiv style)
% Citation keys match the project BibTeX / references.txt

\section{Methodology}
\label{sec:methodology}

This section presents the threat model, data construction procedure, defense designs, and evaluation methodology used in this study.
Our objective is to determine whether lightweight, model-agnostic input sanitization can reduce the success of prompt-injection attacks while preserving acceptable behaviour on benign user requests.
The methodological choices below are guided by three principles.
First, defenses should operate before inference so that they remain portable across target models.
Second, security should be measured both at the filter and at the end-to-end response, because a sanitizer miss does not automatically imply model compliance.
Third, claims of robustness should be supported by controlled stress tests rather than by a single static wording of the test set.

\subsection{Threat Model and System Architecture}
\label{sec:threat-model}

Prompt injection arises when adversarial natural-language content causes a language model to disregard its intended instructions, disclose restricted information, or adopt an attacker-chosen role~\cite{owasp_llm,perez2022ignore,greshake2023indirect,yi2024jailbreak}.
Unlike classical software injection, the malicious payload is expressed in unconstrained language and can be embedded in otherwise fluent requests, which makes purely syntactic defenses incomplete~\cite{11359704,correia2026systematic}.
We focus on \emph{direct} prompt injection, in which the malicious content is contained in the user message itself, rather than introduced indirectly through retrieved documents, web pages, or external tool outputs~\cite{greshake2023indirect,hines2024spotlighting,wu2025epistemic}.
This setting corresponds to the primary risk category identified in the OWASP Top~10 for Large Language Model Applications~\cite{owasp_llm} and remains a standard starting point for defense evaluation before extending to richer agentic environments~\cite{debenedetti2024agentdojo}.

\paragraph{Adversary capabilities.}
We assume an adversary who can submit arbitrary free-form text to a deployed instruction-following model.
The adversary may employ instruction overrides, persona hijacking, multi-step pivots, hypothetical or academic framing, and mild paraphrase, but does not modify model weights, system prompts, or the sanitizer internals.
The defender observes only the incoming user text and must decide whether to allow or refuse it before generation.
No access to privileged system messages beyond what is implied by ordinary chat usage is assumed for the attacker in the direct-injection setting studied here.

\paragraph{Defender capabilities and constraints.}
The defender may inspect the candidate prompt, compute handcrafted or learned risk signals, and block the request.
The defender may not retrain the target model, redesign its decoding strategy, or rely on structured query formats that change the application interface~\cite{struq2024,secalign2024}.
Accordingly, we place a sanitizer as a pre-inference gate in front of the target model.
For each incoming prompt, the sanitizer either forwards the text unchanged or blocks it.
Blocked prompts never reach the model; unblocked prompts are answered in the usual way.
This architecture is consistent with input sanitization and guardrail evaluation practice~\cite{chowdhury2025prepsilonepsilonmptsanitizingsensitiveprompts,geng2025pisanitizer,shi2025promptarmor,zizzo2025adversarial}, and is complementary to training-time or interface-level defenses, which we treat as related but out of scope.

\paragraph{Scope.}
The evaluation corpus is intentionally compact and curated.
The study should therefore be read as a controlled, reproducible comparison of sanitizer designs under a fixed threat model, rather than as a production security certification or an exhaustive assessment of multi-tool agents and retrieval-mediated attacks~\cite{debenedetti2024agentdojo,greshake2023indirect}.
Within this scope we examine defenses of increasing sophistication: an unprotected baseline, a regular-expression filter, a keyword heuristic, a hybrid handcrafted context-aware sanitizer with hard triggers and soft secondary scoring (Method~A), and a learned soft-fusion variant of the same signal family (Method~B).
Lexical filters remain common in practice because they are transparent and inexpensive~\cite{openrouter2026promptinjection,getstream2024prompt,mend2024prompt}, yet they are known to degrade under paraphrase and stylistic rewriting~\cite{yi2024jailbreak,wang2025promptsleuth}.
The context-aware methods therefore incorporate semantic similarity and linguistic naturalness cues, following the broader direction of semantic and context-sensitive detection~\cite{wang2025promptsleuth,kholkar2025capture,bair2025prompt}.

\subsection{Dataset}
\label{sec:dataset}

\subsubsection{Construction principles}

Effective evaluation of input sanitizers requires both clearly malicious prompts and realistic benign requests.
If the benign distribution is unrealistically clean, false-positive rates appear artificially low and usability costs are hidden.
If the attack distribution is dominated by a few stereotyped phrases, exact string matching appears artificially strong and generalizations to rewritten attacks are overstated.
We therefore assembled a mixed English corpus from several public sources, imposed a uniform annotation schema, documented licensing, and reserved a stratified held-out test partition for final reporting.

The corpus is designed to support three uses simultaneously: estimation of unprotected attack prevalence, comparison of lexical and context-aware defenses, and diagnosis of over-blocking on ordinary and technical benign language.
It is not intended to exhaust the space of multilingual, multimodal, or tool-mediated attacks~\cite{yeo2025multimodal,debenedetti2024agentdojo}.

\subsubsection{Corpus composition}

The resulting dataset contains 350 prompts, of which 200 are labeled \emph{injection} and 150 are labeled \emph{benign}.
Each instance consists of an identifier, the prompt text, and a binary label.
This balance intentionally retains a substantial benign share so that false-positive behaviour remains statistically visible on the test partition.

Injection examples are drawn in equal measure from four sources: a community jailbreak collection, a compact jailbreak corpus, a public prompt-injection benchmark, and a malicious-prompt detection dataset.
Benign examples combine public non-malicious prompts with self-authored everyday requests covering ordinary information-seeking, explanation, and writing tasks (Table~\ref{tab:data-sources}).
Taken together, the mixture covers instruction overrides, persona hijacking, policy circumvention framings, and commonplace user behaviour in a single evaluation pool.

\begin{table}[t]
  \centering
  \caption{Composition of the evaluation corpus (\(N=350\)).}
  \label{tab:data-sources}
  \begin{tabular}{lrrlp{4.6cm}}
    \hline
    Slice & \(n\) & Share (\%) & Role & Origin \\
    \hline
    Community jailbreaks & 50 & 14.3 & injection & Public jailbreak prompt collection \\
    Compact jailbreaks & 50 & 14.3 & injection & Compact jailbreak corpus \\
    Prompt-injection benchmark & 50 & 14.3 & injection & Public injection / benign dataset \\
    Malicious-prompt set & 50 & 14.3 & injection & Malicious prompt detection dataset \\
    \textit{Injection subtotal} & \textit{200} & \textit{57.1} & --- & --- \\
    \hline
    Prompt-injection benchmark & 100 & 28.6 & benign & Same public dataset (benign subset) \\
    Self-authored requests & 50 & 14.3 & benign & Everyday benign queries \\
    \textit{Benign subtotal} & \textit{150} & \textit{42.9} & --- & --- \\
    \hline
    \textbf{Total} & \textbf{350} & \textbf{100.0} & --- & --- \\
    \hline
  \end{tabular}
\end{table}

Preprocessing is deliberately light in order to preserve natural wording.
Empty strings are discarded, and all malicious source labels are mapped to the unified injection class.
Random sampling, where required, uses a fixed seed to support reproducibility; one source contributes its first cleaned subset in source order.
Licensing and attribution for each constituent source are recorded separately, and the combined corpus is treated as a non-commercial research artifact.

\subsubsection{Partitioning}

We form stratified training, validation, and test partitions that preserve the injection/benign ratio in each fold (Table~\ref{tab:splits}).
Stratification is important because a chance imbalance on a small test set would otherwise dominate Bypass Rate and false-positive estimates.

The training partition supports three roles: construction of reference embeddings for semantic similarity, calibration of perplexity percentiles, and supervised training of Method~B's fusion model.
The validation partition is used solely to select Method~B's operating threshold, thereby avoiding threshold shopping on the test set.
All primary metrics are reported on the held-out test partition of 53 prompts (30 injection, 23 benign).
Partitions are frozen prior to evaluation and are not reshuffled during experimentation.

\begin{table}[t]
  \centering
  \caption{Stratified dataset partitions.}
  \label{tab:splits}
  \begin{tabular}{lrrrr}
    \hline
    Partition & Total & Injection & Benign & Share \\
    \hline
    Training & 245 & 140 & 105 & 70.0\% \\
    Validation & 52 & 30 & 22 & 14.9\% \\
    Test & 53 & 30 & 23 & 15.1\% \\
    \hline
  \end{tabular}
\end{table}

\subsection{Defense Designs}
\label{sec:defenses}

Each defense maps an input prompt to a binary decision: allow or block.
When a prompt is blocked, no target-model generation is performed for that request.
This common interface allows all methods to be compared under an identical evaluation harness and isolates differences in filtering policy from differences in decoding.

\subsubsection{Unprotected baseline}

The baseline applies no filtering and forwards every prompt to the target model.
It estimates unconstrained attack prevalence under the chosen target and judge configuration, and it serves as the reference against which relative security gains are measured~\cite{zizzo2025adversarial,promptfoo2024owasp}.
Without this condition, reductions in Bypass Rate or True ASR cannot be interpreted as defense effects rather than as intrinsic model refusal.

\subsubsection{Regular-expression filter}

The regular-expression defense maintains a library of 45 case-insensitive patterns targeting recurring attack forms: instruction overrides, persona and role hijacks, named jailbreak personas, unrestricted-mode language, system-prompt exfiltration requests, selected non-English overrides, and common hypothetical or academic framings used to circumvent policy.
A match causes the prompt to be blocked.

Pattern filters are attractive in deployment because they are transparent, auditable, and computationally inexpensive~\cite{owasp_llm,openrouter2026promptinjection}.
Their central limitation is dependence on surface form: an adversary who preserves intent while changing wording can often evade exact patterns.
We therefore treat the regular-expression filter as a strong lexical baseline rather than as a complete solution.

\subsubsection{Keyword heuristic}

The keyword defense assigns suspicion weights to 37 phrases associated with injection and jailbreak language.
Matched weights are accumulated, and the prompt is blocked when the total reaches or exceeds a threshold of 0.5.
By combining multiple weak cues, the heuristic can fire on prompts that contain several mildly suspicious fragments even when no single regular expression matches~\cite{getstream2024prompt,mend2024prompt}.

The same property creates a usability risk: technical benign language that reuses words such as ``override'', ``kill'', or ``jailbreak'' in harmless senses may accumulate enough weight to be blocked.
The keyword method therefore illustrates both the appeal and the brittleness of lexicon-driven moderation.

\subsubsection{Method A: Handcrafted context-aware sanitization}
\label{sec:method-a}

Method~A is the primary handcrafted defense proposed in this work.
We frame it explicitly as a \emph{hybrid} architecture rather than as a pure ensemble scorer.
A high-confidence hard-trigger layer provides the main protection when semantic or structural evidence is decisive, while a weighted multi-signal score supplies secondary coverage for borderline cases that do not meet those veto conditions.
Additional heuristic signals improve robustness and interpretability even when their standalone contribution is limited.
The design rests on a practical observation repeatedly noted in the jailbreak and injection literature: paraphrased and role-play attacks often evade exact string matching while still leaving semantic, structural, or stylistic traces~\cite{wang2025promptsleuth,kholkar2025capture,yi2024jailbreak,bair2025prompt}.

\paragraph{Design rationale.}
Lexical detectors contribute precision on known attack phrases, but they miss intent-preserving rewrites.
Semantic similarity contributes coverage across wording changes by comparing the candidate prompt to known injections in embedding space~\cite{reimers2019sentence,wang2020minilm}.
Structural detectors contribute sensitivity to the rhetorical shape of an injection, such as persona assignment or multi-step pivots.
Perplexity contributes a weak naturalness prior that can raise suspicion for highly unnatural strings without being trusted in isolation~\cite{radford2019language}.
In the hybrid design, strong semantic evidence is elevated to a primary veto, while the remaining cues participate mainly through the complementary soft layer.

\paragraph{Signal inventory.}
The eight signals are as follows.
\begin{enumerate}
  \item \emph{Regular expression:} whether the pattern filter would block the prompt, reused as one vote among several rather than as the sole decision maker.
  \item \emph{Keyword:} the keyword suspicion score, capped so that a long list of weak matches cannot dominate the aggregate.
  \item \emph{Semantic similarity:} cosine similarity between a compact sentence embedding of the candidate prompt and embeddings of training injection prompts.
    The maximum similarity is mapped to a stepped risk score so that strong resemblance to known attacks contributes more than marginal overlap.
  \item \emph{Intent:} structural patterns that combine instruction override with disclosure or policy-break requests, persona assignment, conversation resets, and named jailbreak modes.
  \item \emph{Role-play:} cues associated with persona hijacking and unrestricted ``do anything'' style characters that package harmful requests inside fiction.
  \item \emph{Instruction shift:} multi-step pivots in which a benign opening is followed by a conflicting directive.
  \item \emph{Objective conflict:} juxtaposition of an apparently legitimate task with an incompatible malicious goal in the same prompt.
  \item \emph{Perplexity:} relative linguistic unnaturalness under a small causal language model, calibrated on the training partition and treated as a supporting cue rather than a standalone detector.
\end{enumerate}

\paragraph{Primary layer: hard triggers.}
Method~A first applies three hard vetoes that block immediately when
(i)~the stepped semantic score exceeds 0.6,
(ii)~the intent detector is active, or
(iii)~role-play is active and the keyword score exceeds 0.3.
Among these, the semantic hard trigger is intended as the principal high-confidence defense; the other vetoes provide additional structural coverage.
Perplexity never triggers a hard veto, because elevated perplexity also appears in benign creative writing, code-like instructions, and terse technical queries.

\paragraph{Secondary layer: soft aggregation.}
If no hard trigger fires, Method~A falls back to a weighted soft score over the eight signals.
The weights are
0.12 (regular expression),
0.12 (keyword),
0.20 (semantic similarity),
0.18 (intent),
0.10 (role-play),
0.08 (instruction shift),
0.08 (objective conflict),
and 0.12 (perplexity).
A prompt is soft-blocked when the aggregated score reaches or exceeds 0.40.
This secondary layer is meant to handle borderline multi-signal cases that do not satisfy hard-trigger conditions, rather than to serve as the sole source of protection.
The final decision blocks if either a hard veto fires or the soft threshold is crossed.

\paragraph{Computational placement.}
Sentence embedding and perplexity components are executed on CPU by design.
This choice keeps accelerator memory available for the target model and allows sanitizer latency to be measured as an independent pre-inference cost.

\subsubsection{Method B: Learned context-aware sanitization}
\label{sec:method-b}

Method~B addresses a natural follow-up question: if Method~A's signals are useful, can supervised fusion replace hand-chosen weights and hard rules?
Method~B therefore retains the same signal family, expands the semantic representation, and replaces rule-based aggregation with probabilistic classification.
The approach is related in spirit to learned detectors and semantic intent modeling~\cite{wang2025promptsleuth,shi2025promptarmor}, while remaining a lightweight classical model rather than a large dedicated guardrail language model.

\paragraph{Feature representation.}
Instead of one stepped semantic score, Method~B uses three continuous semantic features: the maximum similarity to training injections, the mean of the top-three injection similarities, and a benign-contrast term.
The contrast term reduces risk when a prompt also resembles ordinary benign text, which is intended to limit false positives on near-boundary benign requests.
Together with regular-expression, keyword, intent, role-play, instruction-shift, objective-conflict, and perplexity signals, these features form a ten-dimensional representation.

\paragraph{Learning and calibration.}
Features are standardized and passed to a regularized logistic regression classifier with balanced class weighting, trained on the training partition.
The operating threshold is selected on the validation partition by maximizing F1 score over a dense grid of candidates; ties are resolved in favour of lower Bypass Rate and then lower false-positive rate.
At test time, a prompt is blocked when the predicted injection probability meets or exceeds this threshold.

\paragraph{Relation to Method A.}
Unlike Method~A's hybrid design, Method~B employs no hard vetoes and therefore relies entirely on soft probabilistic fusion.
Any gain or loss relative to Method~A is attributable to feature expansion and learned calibration rather than to an explicit high-confidence override layer.
This contrast is central to later ablation analysis: it asks whether strong empirical security requires hard semantic gating, as in Method~A, or whether learned soft decisions alone suffice.

\subsubsection{Optional lexical classifier}

For reference, we also consider a classical bag-of-\(n\)-grams logistic regressor with TF--IDF features~\cite{salton1988term}.
This classifier provides a supervised lexical baseline that can exploit co-occurrence patterns beyond hand-written dictionaries.
It is retained as an optional reference point and is not part of the primary defense comparison unless separately reported.

\subsection{Two-Layer Evaluation}
\label{sec:evaluation}

Sanitizer evaluation is easy to mis-specify.
A method that blocks nearly everything will look excellent on Bypass Rate while destroying usability.
A method that misses many attacks may still appear acceptable if the target model refuses on its own.
We therefore evaluate security at two layers and always report usability alongside attack metrics, following contemporary guardrail benchmarking practice~\cite{zizzo2025adversarial,promptfoo2024owasp}.

\paragraph{Target model.}
Unblocked prompts are answered by an instruction-tuned open model, Qwen2.5-3B-Instruct, under greedy decoding.
Using a single open target keeps the comparison focused on sanitizer behaviour rather than on opaque differences among proprietary backends.
Greedy decoding removes sampling variance from the principal security tables.

\paragraph{Compliance judgment.}
Layer~2 asks whether an unblocked injection actually succeeded.
For each such case, a separate judge from the same model family receives the original prompt and the target response and labels the outcome as complied or refused~\cite{zizzo2025adversarial,yi2024jailbreak}.
Judgments are obtained from an explicit final label; unparseable outputs are treated conservatively as refusals.
Long inputs are truncated before judging, and the judge is checked on a small hand-authored calibration set before the main evaluation.

\paragraph{Definition of a successful attack.}
An attack is counted as successful only when three conditions hold simultaneously: the example is labeled injection, the sanitizer fails to block it, and the judge deems the model response compliant.
This definition ensures that end-to-end risk is not equated with a filter miss alone, and that model refusals after a bypass are credited correctly.

\paragraph{Primary metrics.}
All primary security and usability figures in this paper are computed on the frozen held-out test partition of 53 prompts (30 injection, 23 benign) under the full two-layer pipeline (sanitizer, target model, and judge).
We report four quantities:
\begin{itemize}
  \item \emph{Bypass Rate:} the fraction of injection prompts left unblocked (Layer~1);
  \item \emph{Detection Rate:} the complementary fraction of injection prompts blocked;
  \item \emph{False Positive Rate (FPR):} the fraction of benign prompts incorrectly blocked;
  \item \emph{True Attack Success Rate (True ASR):} the fraction of injection prompts that both bypass the sanitizer and elicit a complied response (Layer~2).
\end{itemize}
These two security metrics are deliberately kept distinct.
Bypass Rate measures only whether the sanitizer failed to block an injection.
True ASR additionally requires judged model compliance.
We therefore never collapse them into a single undifferentiated ``ASR'' figure: an unblocked injection that the target refuses contributes to Bypass Rate but not to True ASR.
Confusion matrices are computed at Layer~1 only, using injection versus blocked decisions, so that judge labels do not enter those counts.

\subsection{Robustness Protocols}
\label{sec:robustness}

A defense that performs well only on the original wording of a small test set is easy to overstate.
We therefore include three complementary robustness protocols aligned with common red-teaming practice~\cite{promptfoo2024owasp,zizzo2025adversarial}.

\paragraph{Paraphrase robustness.}
Mild synonym substitutions are applied to a subset of training injection prompts (for example, replacing ``ignore'' with ``disregard'' or ``act as'' with ``behave as'').
For each defense we compare Bypass Rate on original and paraphrased text under sanitizer-only evaluation.
Because no target-model call is involved, the measurement isolates lexical brittleness of the filter under controlled wording drift.
The paraphrase operator is intentionally simple; its purpose is diagnostic rather than adversarial optimality.

\paragraph{Edge-case benign evaluation.}
Security metrics alone hide a different failure mode: over-blocking.
We evaluate a fixed collection of 80 technical but benign prompts that reuse attack-like vocabulary in harmless contexts, such as process control, style overrides, or device-related uses of ``jailbreak''.
The fraction blocked by each defense estimates false-positive risk beyond the standard benign test partition.
High edge-case false-positive rates indicate that a method may be difficult to deploy in developer-facing or systems-oriented assistants~\cite{bair2025prompt,getstream2024prompt}.

\paragraph{Adaptive rewriting.}
Static test prompts can understate attacker effort even when paraphrase is mild.
Each test injection is therefore rewritten into five stylistic variants---role-play, research framing, hypothetical framing, indirect phrasing, and multistep instruction---using a lightweight sequence-to-sequence model.
The resulting adaptive pool contains up to 150 prompts and is evaluated with the full two-layer pipeline.
This protocol stresses defenses under controlled style shift.
It is not intended as a substitute for strong optimization-based jailbreaks, nor for agentic indirect-injection benchmarks~\cite{debenedetti2024agentdojo,greshake2023indirect,hines2024spotlighting}; rather, it asks whether a defense still functions when the same malicious goals are restated.

\subsection{Ablation Analysis}
\label{sec:ablation}

The context-aware defenses contain several interacting ingredients, so aggregate accuracy alone does not reveal which components are responsible for observed gains.
We therefore conduct leave-one-component-out ablations on the same held-out test partition used for primary reporting.

For Method~A, we study two complementary families that test the hybrid claim.
The first disables hard vetoes and removes each soft signal in turn, thereby measuring secondary soft-layer coverage when the primary hard-trigger layer is unavailable.
The second compares the full system with a soft-only variant and with variants that disable each hard veto individually, thereby isolating whether headline security depends on hard triggers---especially the semantic veto---rather than on soft aggregation.
For Method~B, selected feature groups are zeroed at inference time---including an all-semantic condition---while the trained fusion model remains fixed.
This distinguishes whether learned performance depends on semantic features specifically or on the broader feature set.

Each variant is summarized by Bypass Rate, True ASR, false-positive rate, and percentage-point changes relative to the corresponding full configuration.
These ablations underwrite later claims about the relative importance of hard semantic gating, soft fusion, and individual signal families.

\subsection{Efficiency Considerations}
\label{sec:runtime}

A pre-inference gate that is accurate but slow may still be impractical in interactive applications.
We therefore treat latency as part of the methodological evaluation rather than as an afterthought.
Wall-clock time is measured for semantic scoring, perplexity scoring, full Method~A sanitization, and the lexical baselines under a fixed prompt sample and warm-up protocol.
Because auxiliary embedding and language-model components run on CPU, the resulting figures characterize sanitizer overhead independently of target-model GPU inference and can be interpreted as the added cost of inserting the gate into a serving path.
